\documentclass[10pt,a4paper]{amsart}
\usepackage[margin=19mm]{geometry}
\usepackage{amsmath,amssymb,mathtools}
\usepackage[scr=rsfs,cal=euler]{mathalfa}
\usepackage{xfrac}
\usepackage[unicode,hidelinks,bookmarks=false]{hyperref}
\newcommand{\ie}{i.e.\ }
\newcommand{\cf}{cf.\ }
\newcommand{\resp}{respectively\ }
\newcommand{\Subsection}{\subsection}
\providecommand{\nobreakdash}{\nobreak\hbox{-}}
\setlength{\emergencystretch}{2em}
\multlinegap0pt
\usepackage{amssymb}
\usepackage{xcolor}
\newtheorem{lem}{Lemma}
\newtheorem{prop}{Proposition}
\title{Nonexistence of Approximately Dual Frames via Weighted Composition Operators}
\author{Trevor Camper, Dongwei Chen, Shuang Guan}
\date{Source: arXiv:2609.12357v1 (CC BY 4.0)}
\begin{document}
\maketitle

\noindent\emph{Source and licence.} Nonexistence of Approximately Dual Frames via Weighted Composition Operators. Version arXiv:2609.12357v1 (CC BY 4.0), distributed by arXiv under the Creative Commons Attribution 4.0 International licence, http://creativecommons.org/licenses/by/4.0/, as recorded in the arXiv licence metadata for this exact version and shown on the abstract page https://arxiv.org/abs/2609.12357v1. Full licence notice: \texttt{licenses/arxiv-2609.12357-CC-BY-4.0.txt}.\par\smallskip

\noindent\textit{Editorial scope.} Counted unit: the article's prop fock\_main (source0.tex lines 716--727). The article's complete original proof of it is delivered with it (source0.tex lines 728--776, one \texttt{proof} environment of 49 lines). No mathematics is written, repaired or re-derived: the statement and the proof are the article's own lines, in the article's own notation, and no retained line is substituted or re-flowed.  Retained uncounted as the source-local context the proof needs: lines 641--654.  Nothing was omitted from any retained span, so the package carries no omission record.  No retained line contains a citation, so the wrapper carries no bibliography and no bibliography entry.  No retained line contains a figure environment, an image-inclusion command or drawing code, so no image is delivered and the package declares no asset dependency.  Editorial formatting only, in addition: an AMS \texttt{amsart} preamble that restates the article's own declarations and macros used by the retained lines, namely the environments \texttt{lem}, \texttt{prop} on separate counters, together with the standard packages those lines need.  The retained environments print in this document as Lemma 1, Proposition 1 in order of appearance, whereas the article prints the same items with the numbering of its own full document.  The complete native source of the article as distributed in the arXiv e-print for this exact version is bundled unchanged as \texttt{sources/210108/source0.tex}.
\begin{lem}\label{wco_fock_lem}
    Let $W = W_{\psi,\varphi}$ be a bounded, invertible weighted
composition operator on $\mathcal{F}^2(\mathbb{C}^n)$. Then $W$ is a scalar
multiple of a unitary operator. More precisely, for some $A$ unitary and $b \in \mathbb{C}^n$,
\[
\varphi(z) = Az + b, \qquad
\psi(z) = \psi(0)\, e^{-\langle z, A^*b\rangle},
\]
and $W$ is a constant multiple of an isometry: for every
$f \in \mathcal{F}^2(\mathbb{C}^n)$,
\[
\|Wf\| = r\|f\|,\; \text{where} \; r := |\psi(0)|\, e^{|b|^2/2}.
\]
\end{lem}

\begin{prop}\label{fock_main}
Let $\{f_n\}$ be a frame for $\mathcal{F}^2(\mathbb{C}^n)$ with optimal frame
bounds $0 < A \le B < \infty$, and let $0 \le \varepsilon < 1$. Then
there exists a bounded weighted composition operator $W$ on
$\mathcal{F}^2(\mathbb{C}^n)$ such that $\{Wf_n\}$ is an
$\varepsilon$-approximate dual frame to $\{f_n\}$ if and only if
\[
\frac{B - A}{A + B} \le \varepsilon,
\qquad\text{equivalently,}\qquad
\frac{B}{A} \le \frac{1 + \varepsilon}{1 - \varepsilon}.
\]
\end{prop}

\begin{proof}
The two inequalities are equivalent since
$B(1-\varepsilon) \le A(1+\varepsilon)$ if and only if
$B - A \le \varepsilon(A + B)$.

($\Leftarrow$) Let $c = \tfrac{2}{A+B}$ and take $W = cI$, a weighted
composition operator with constant weight $\psi \equiv c$ and
$\varphi = \mathrm{Id}$, where $\mathrm{Id}$ is the identity map for $\mathbb{C}^n$. Let $S$ be the frame operator of $\{f_n\}$.  Since $S$ is self-adjoint with
$\sigma(S) \subseteq [A, B]$, it holds that
\[
\|I - WS\| = \|I - cS\|
= \max_{t \in \sigma(S)} |1 - ct|
\le \max_{t \in [A,B]} |1 - ct|
= \frac{B-A}{A+B} \le \varepsilon.
\]
Note
that $\{Wf_n\} = \{cf_n\}$ is a frame trivially. Therefore, $\{Wf_n\}$ is an $\varepsilon$-approximate dual frame to $\{f_n\}$.

($\Rightarrow$) Suppose $\|I - WS\| \le \varepsilon < 1$. Then $WS$
is invertible by a Neumann series argument, and since $S$ is
invertible, so is $W$. By Lemma \ref{wco_fock_lem} there is
$r > 0$ with $\|Wg\| = r\|g\|$ for all $g \in \mathcal{F}^2(\mathbb{C}^n)$.
Then for every unit vector $f$,
\[
\bigl|\,1 - r\|Sf\|\,\bigr| = \bigl|\,\|f\| - \|WSf\|\,\bigr|
   \le \|(I-WS)f\| \le \varepsilon,
\]
and
\begin{equation}\label{eq:4.3ineq}
1 - \varepsilon \le r\|Sf\| \le 1 + \varepsilon,
\qquad \|f\| = 1.
\end{equation}
Since the frame bounds are optimal and $S$ is positive and self-adjoint,
$\|Sf\|^2 = \langle S^2 f, f\rangle$ which indicates that
\[
\inf_{\|f\|=1} \|Sf\| = A
\qquad\text{and}\qquad
\sup_{\|f\|=1} \|Sf\| = B
\]
both attained in the limit along the approximate eigenvectors of $S$ at
$A$ and $B$ respectively. Taking such approximate eigenvectors in
\eqref{eq:4.3ineq} yields
\[
rA \ge 1 - \varepsilon
\qquad\text{and}\qquad
rB \le 1 + \varepsilon,
\]
and dividing gives $B/A \le (1+\varepsilon)/(1-\varepsilon)$.
\end{proof}

\end{document}
